\documentclass[12pt]{article}

\usepackage[a4paper, total={6in, 8in}]{geometry}
\usepackage{textcomp}

\begin{document}
\setlength{\parindent}{0pt}

\def \t {\textrightarrow}
\def \v {\vspace{1em}}
\def \bi {\begin{itemize}}
\def \ei {\end{itemize}}
\def \s[#1] {\section*{#1}}
\def \ss[#1] {\subsection*{#1}}
\def \sss[#1] {\subsubsection*{#1}}

\s[Carducci - Stazione di autunno]
In "Inno a satana" è presente la figura del treno, ma totalmente diversa \t qua è un mostro empio \\
Il giovane Carducci ha visto nel treno l'emblema del progresso \t e solo in quello + matura c'è visione + pericolosa del treno \t delusione per questa macchina, che divide \\
Viso pallido, leggermente roseo \t candida, pura \\
Sole di giugno, giovane perchè è l'inizio dell'estate \t crine termine elettivo, che si trova nell'epica \t classicismo di Carducci emerge \\
Riflessi castanei del crine \t anastrofe \t verso 45, evoluzione dell'immagine della donna e femminilità \t sensualità rispettosa della tradizione per dolcezza \\
La molle guancia \t nel senso irrorata di sudore \\
Aureola \t i miei sogni soggetto, cingono la persona \\
Il suo essere l'amante di Carducci, non implica che siano moralmente cattive persone \\
Il dolore della morte del figlio ha creato frattura tra i due \t lacerazione dell'animo ha forse allontanato i due \\
Verso 39 riprende ripiegamento interiore \t caligine, sentore boudeleriano \t ebbro, battello ebbro \t senza di lei perde consapevolezza di se \\
Esclamazione del verso 51-52-52 \t eco pascoliano \t pioggia di foglie gelida, continua e muta \t foglie sono emblema dell'autunno dell'allontanamento \\
Verso 55 eterno, climax ascendente e chiasmo \t che che \t è il dolore universale, tutti gli uomini sperimentano la privazione e l'abbandono \\
Padre privato del figlio \t non ha identificazione lessicale \t un genitore che perde un figlio è innaturale \\
Io voglio \t pioggia, caligine = univeso boudeleriano \t tedio è stato di sopsensione che non permette di calarsi ne nella gioia ne nel dolore \\
Tedio descrive stato di sospensione \\
Testo unisce elementi personali, realistici, individualistici, che unisce il passato letterario al presente letterario \\
Immagine di un dolore universale e individuale allo stesso tempo \t un canto corale e individuale \\
Inno a satana, vedere porzione dove compaioni questi \t dove carducci giovane inneggia al progresso \t definito il treno bello e orribile

\v

D'annunzio, inizio inettitudine, pascoli + intimistica e rispettosa del proprio dolore, il decadentismo che apre sempre verso l'europa, e un estetismo estremo con Wilde a Hyusmann con A Reboir = complessita del decadentismo \\
Apre le porte al nuovo romanzo e al cinema attuale \\
Romanzo moderno ha diversi indicatori, ma accumanto dalla crisi dell'io e dell'identità \\
Ogni pelicola ha finale sospeso \t siamo noi a dare un finale, una conclusione \t l'eco di manzoni è lontano \\
Si parla di una nuova era, e di diversi volti che si incastonano in un prisma \t a seconda della rifrazione della luce, si vedono un cangiare diverso dei colori \\
La fase che accompagna la fine del 900 è caratterizzata da personalita con denominatore comune: contrastare manzoni \\
Questo si configura in modi diversi \t Tozzi parla dell'inettitudine \t trateggia la figura dell'inetto, in un orizzonte circoscritto geograficamente \\
Questo sara il punto di partenza per Pirandello etc. \\
Scalpitano elementi culturali \t Compte nascono scienze umane \t uomo con le domande di senso \t si ricorda la nascita della fisionomia \t \\
Uomo si sente ingabbiato per le prime ravvisaglie di un cambiamento epocale \t uomo è completamente perso \\
Fogazzaro nella figura di Malombra presente figura femminile che vive con la consapevolezza del presente e i fantasmi del passato \t e analizza lo scontro tra chiesa e riformismo \t la scienza e la fede da equilibrare in quale rapporto? \\
Parlare di modernismo significa riflettere sulla visione progressista \\
Visione del tempo differente se affidata alla consistenza dei fatti \t Bergson lo sa bene \\
In Tozzi si mosra la destrutturazione di un eroe, è l'inetto, che in Joyce sarà l'ulisse \t egli vive un odissea in un giorno \\
Il padre dello stream of consciusness è James Joyce

\s[Pirandello]
Vita e arte sono fortemente collegate, la pazzia è un filo conduttore \t sua moglie impazzisce \t antonietta portulano, a causa di uno scoppio di una zolfatara \\
Il patrimonio maturale si anniente, e lo schock della moglie porta a questo \t evento coinvolge anche Pirandello, che vive questo \\
La zolfatara era di proprieta del suocero (e quindi della moglie) \t sempre in sicilia, un personaggio che configura gli ultimi sprazzi del 900 \t muore nel 36 \\
Pirandello è personalità inimitabile \t squarto culturale, imprescindibile per capire il presente \t pirandello descrive l'evolversi del romanzo in sfaccettature che si muovono gia da "L'esclusa" \\
La triade dei romanzi è l'esclusa, il fu mattia pascal e uno nessuno centomila \\
Lui pero è anche drammaturgo, e lascia l'evoluzione del teatro \t metateatro che parla di se stesso, che pone i riflettoeri sul dualismo tra personaggio e persona \\
1906-7 fa scritti teorici, soprattutto sull'umorismo \t anziana signora truccata come un pappagallo, sono le nostre idee precostituite \t nasce il riso, che sorge dall'avvertimento del contrario \\
Dopo il riso naturale, sorge una riflessione innaturale \t e sta dietro un dolore, una volonta di afferrare il tempo, e di destare scandalo
\end{document}
